% GENERATED FILE. Edit canonical lessons, then run npm run generate:books.
% canonical-source-hash: fnv1a64:4f7ee7e211bf9969
% canonical-lessons: TE-C104-tiram, TE-C104-karyalayam, TE-C104-tota, TE-C104-mula, TE-C104-malupu

\chapter{Shore, Office, Garden, Corner, Turn}
\label{ch:te-shore-office-garden-corner-turn}
\hlchaptermodality{104}

\begin{chapteropening}
\textbf{By the end of this chapter:} \emph{I can say a shore, an office, a garden, a corner and a turn.}

Complete the last lesson of chapter 104: i can say a shore, an office, a garden, a corner and a turn.
\end{chapteropening}

\section[tīraṁ]{\te{తీరం} (tīraṁ) — a shore, a coast}

\label{lesson:TE-C104-tiram}



\begin{quote}
Before the new one: say the Telugu for a bank, then the Telugu for a market.
\end{quote}

\subsection*{You'll want to know: \te{తీరం}}
\textbf{\te{తీరం}} — \emph{tīraṁ} — \textquotedblleft{}a shore, a coast\textquotedblright{}.

\begin{grammarlens}[title={What you've built}]
One more word for this chapter.
\end{grammarlens}

\subsection*{Your turn}
\begin{itemize}
\raggedright
  \item \emph{Say it:} \emph{tīraṁ}
  \item \emph{Say it:} \emph{tīraṁ}, once more
  \item \emph{Say it:} say \emph{tīraṁ}
  \item \emph{Recall:} say the Telugu for a bank, then the Telugu for a market, then say \emph{tīraṁ} again
\end{itemize}

\begin{tcolorbox}[breakable,colback=teal!4,colframe=teal!35!black,title={Before you move on}]
What does \te{తీరం} mean? (\textquotedblleft{}A shore, a coast\textquotedblright{}.) Say it once more.
\end{tcolorbox}
\section[kāryālayaṁ]{\te{కార్యాలయం} (kāryālayaṁ) — an office}

\label{lesson:TE-C104-karyalayam}



\begin{quote}
Before the new one: say the Telugu for a market, then the Telugu for a shore.
\end{quote}

\subsection*{You'll want to know: \te{కార్యాలయం}}
\textbf{\te{కార్యాలయం}} — \emph{kāryālayaṁ} — \textquotedblleft{}an office\textquotedblright{}.

\begin{grammarlens}[title={What you've built}]
One more word for this chapter.
\end{grammarlens}

\subsection*{Your turn}
\begin{itemize}
\raggedright
  \item \emph{Say it:} \emph{kāryālayaṁ}
  \item \emph{Say it:} \emph{kāryālayaṁ}, once more
  \item \emph{Say it:} say \emph{kāryālayaṁ}
  \item \emph{Recall:} say the Telugu for a market, then the Telugu for a shore, then say \emph{kāryālayaṁ} again
\end{itemize}

\begin{tcolorbox}[breakable,colback=teal!4,colframe=teal!35!black,title={Before you move on}]
What does \te{కార్యాలయం} mean? (\textquotedblleft{}An office\textquotedblright{}.) Say it once more.
\end{tcolorbox}
\section[tōṭa]{\te{తోట} (tōṭa) — a garden}

\label{lesson:TE-C104-tota}



\begin{quote}
Before the new one: say the Telugu for a shore, then the Telugu for an office.
\end{quote}

\subsection*{You'll want to know: \te{తోట}}
\textbf{\te{తోట}} — \emph{tōṭa} — \textquotedblleft{}a garden\textquotedblright{}.

\begin{grammarlens}[title={What you've built}]
One more word for this chapter.
\end{grammarlens}

\subsection*{Your turn}
\begin{itemize}
\raggedright
  \item \emph{Say it:} \emph{tōṭa}
  \item \emph{Say it:} \emph{tōṭa}, once more
  \item \emph{Say it:} say \emph{tōṭa}
  \item \emph{Recall:} say the Telugu for a shore, then the Telugu for an office, then say \emph{tōṭa} again
\end{itemize}

\begin{tcolorbox}[breakable,colback=teal!4,colframe=teal!35!black,title={Before you move on}]
What does \te{తోట} mean? (\textquotedblleft{}A garden\textquotedblright{}.) Say it once more.
\end{tcolorbox}
\section[mūla]{\te{మూల} (mūla) — a corner}

\label{lesson:TE-C104-mula}



\begin{quote}
Before the new one: say the Telugu for an office, then the Telugu for a garden.
\end{quote}

\subsection*{You'll want to know: \te{మూల}}
\textbf{\te{మూల}} — \emph{mūla} — \textquotedblleft{}a corner\textquotedblright{}.

\begin{grammarlens}[title={What you've built}]
One more word for this chapter.
\end{grammarlens}

\subsection*{Your turn}
\begin{itemize}
\raggedright
  \item \emph{Say it:} \emph{mūla}
  \item \emph{Say it:} \emph{mūla}, once more
  \item \emph{Say it:} say \emph{mūla}
  \item \emph{Recall:} say the Telugu for an office, then the Telugu for a garden, then say \emph{mūla} again
\end{itemize}

\begin{tcolorbox}[breakable,colback=teal!4,colframe=teal!35!black,title={Before you move on}]
What does \te{మూల} mean? (\textquotedblleft{}A corner\textquotedblright{}.) Say it once more.
\end{tcolorbox}
\section[malupu]{\te{మలుపు} (malupu) — a turn, a bend}

\label{lesson:TE-C104-malupu}



\begin{quote}
Before the new one: say the Telugu for a garden, then the Telugu for a corner.
\end{quote}

\subsection*{You'll want to know: \te{మలుపు}}
\textbf{\te{మలుపు}} — \emph{malupu} — \textquotedblleft{}a turn, a bend\textquotedblright{}.

\begin{grammarlens}[title={What you've built}]
One more word for this chapter.
\end{grammarlens}

\subsection*{Your turn}
\begin{itemize}
\raggedright
  \item \emph{Say it:} \emph{malupu}
  \item \emph{Say it:} \emph{malupu}, once more
  \item \emph{Say it:} say \emph{malupu}
  \item \emph{Recall:} say the Telugu for a garden, then the Telugu for a corner, then say \emph{malupu} again
\end{itemize}

\begin{tcolorbox}[breakable,colback=teal!4,colframe=teal!35!black,title={Before you move on}]
What does \te{మలుపు} mean? (\textquotedblleft{}A turn, a bend\textquotedblright{}.) Say it once more.
\end{tcolorbox}
